\skripsifrontchapter{ABSTRACT}

\begin{otherlanguage}{english}
\noindent\textbf{\namamahasiswa, \judulskripsiinggris}\par
\noindent\textbf{Supervisors: \pembimbingsatu{} and \pembimbingdua}\par
\vspace{\baselineskip}

This sample document demonstrates a LaTeX thesis template for organizing preliminary pages, six main chapters, a bibliography, and an appendix. Its purpose is to provide an editable starting point and help users understand the relationship between source structure and document appearance. The included material explains formatting and does not report an empirical study. Student and supervisor identities use generic values, while approval statements, declarations of originality, dates, and signatures require completion according to actual circumstances. The introduction illustrates connections among context, questions, objectives, benefits, and scope. The background chapter provides examples of citations to official documentation and definitions of mathematical symbols. The methodology chapter presents relationships between stages and materials through a table and a vector diagram created within the document. The results chapter demonstrates a compact table, a multipage table with repeated column headings, and cross references to a figure and equations. Its tables contain component names and output descriptions without invented scores, sample sizes, or successful experiments. The discussion chapter explains relationships among evidence, interpretation, and limitations, while the closing chapter illustrates conclusions and recommendations that users must adapt. The bibliography contains genuine online documentation to demonstrate citation management, and the appendix organizes supporting notes. Users must replace every sample passage with accountable research content, check the applicable writing requirements, and review the generated document after the final changes. Successful compilation does not replace examination of content, formatting, references, or administrative completeness before the document is used as an academic submission.

\noindent\textbf{Keywords:} LaTeX, thesis template, document structure, citation, cross reference, format review.
\end{otherlanguage}
